\documentclass[a4paper]{article}

\usepackage[fontset=fandol]{ctex}
\usepackage{amsmath, amssymb}
\usepackage{graphicx}
\usepackage[margin=2.5cm]{geometry}
\usepackage[most]{tcolorbox}
\usepackage{etoolbox}
\usepackage{listings}
\usepackage{hyperref}
\usepackage{booktabs}
\usepackage{subcaption}
\usepackage{float}
\usepackage{tikz}
\IfFileExists{pgfplots.sty}{
    \usepackage{pgfplots}
    \pgfplotsset{compat=1.18}
}{}

\newtcolorbox{knowledgebox}[1]{
    enhanced, colback=blue!5!white, colframe=blue!75!black, colbacktitle=blue!75!black,
    coltitle=white, fonttitle=\bfseries, title=#1, attach boxed title to top left={yshift=-2mm, xshift=2mm},
    boxrule=1pt, sharp corners
}

\newtcolorbox{importantbox}[1]{
    enhanced, colback=yellow!10!white, colframe=yellow!80!black, colbacktitle=yellow!80!black,
    coltitle=black, fonttitle=\bfseries, title=#1, sharp corners
}

\newtcolorbox{warningbox}[1]{
    enhanced, colback=red!5!white, colframe=red!75!black, colbacktitle=red!75!black,
    coltitle=white, fonttitle=\bfseries, title=#1, sharp corners
}

\lstset{
    language=Python,
    basicstyle=\ttfamily\small,
    keywordstyle=\color{blue},
    stringstyle=\color{red!60!black},
    commentstyle=\color{green!60!black},
    breaklines=true,
    frame=single,
    numbers=left,
    numberstyle=\tiny\color{gray},
    captionpos=b,
    extendedchars=false
}

\newcommand{\notetitle}{CS224R 第 5 讲：Off-Policy Actor-Critic}
\newcommand{\noteauthors}{五道口纳什 \& Codex}
\newcommand{\notedate}{2025-04-16}
\newcommand{\videochannel}{Stanford CS224R}
\newcommand{\videopublishdate}{2025-04-16}
\newcommand{\videoduration}{01:09:17}
\newcommand{\videourl}{https://www.youtube.com/watch?v=cRGKc-nAWho}
\newcommand{\videocoverpath}{cover.jpg}

\begin{document}

\begin{titlepage}
\centering
{\Large 课程笔记\par}
\vspace{1.2cm}
{\huge\bfseries \notetitle\par}
\vspace{0.8cm}
{\large \noteauthors\par}
\vspace{0.3cm}
{\large \notedate\par}
\vspace{1.2cm}

\includegraphics[width=0.82\textwidth,height=0.45\textheight,keepaspectratio]{\videocoverpath}\par

\vfill
\begin{tcolorbox}[width=0.9\textwidth, colback=black!2!white, colframe=black!60, sharp corners]
\textbf{视频作者/频道}：\videochannel\par
\textbf{发布日期}：\videopublishdate\par
\textbf{视频时长}：\videoduration\par
\textbf{视频链接}：\href{\videourl}{\nolinkurl{\videourl}}
\end{tcolorbox}
\end{titlepage}

\tableofcontents
\newpage

\section{为什么要更 off-policy：从一次更新走向多次更新}
上一讲的 actor-critic 已经把“估值”和“改策略”拆开，但整体仍偏 on-policy：数据来自当前策略，通常采一批数据只更新一次。这个代价在真实 RL 中非常高，因为与环境交互往往比网络反向传播昂贵得多。

本讲的第一步扩展，就是在同一批由旧策略 \(\pi_\theta\) 收集的数据上，对新策略 \(\pi_{\theta'}\) 做多次梯度更新。讲者提醒，这里的关键不在“多做几次优化”，而在于要正确处理新旧策略分布不一致的问题。

\subsection{importance ratio 的角色}
为了在旧数据上近似新策略目标，需要使用重要性权重
\[
\rho_t=\frac{\pi_{\theta'}(a_t\mid s_t)}{\pi_\theta(a_t\mid s_t)}.
\]
于是 surrogate objective 或等价梯度中，都会出现 \(\rho_t \hat A_t\) 这一项。它的含义是：如果新策略更偏爱这个旧动作，就放大其影响；如果新策略更不偏爱它，就缩小其影响。

这个修正是 off-policy actor-critic 的第一道门槛。没有它，多次更新虽然表面上节省样本，但优化目标已经和真实策略目标脱节。

\subsection{为什么只靠 importance sampling 不够}
讲者也反复强调，importance sampling 不是银弹。若 \(\pi_{\theta'}\) 与 \(\pi_\theta\) 差得太远，则概率比会爆炸或塌缩，带来巨大方差。于是 practical RL 必须想办法约束“每一次策略更新不要走太远”。

\subsection{本章小结}
更 off-policy 的第一步不是 replay buffer，而是先允许一批数据支持多次更新；而一旦这样做，重要性采样与步长约束就成为必须处理的稳定性问题。

\section{PPO：把 actor-critic 做成稳定可用的工程算法}
PPO 是本讲最重要的 practical algorithm。它本质上仍是 actor-critic，但在目标函数和优势估计上加入了若干技巧，使其在样本效率与训练稳定性之间取得实用平衡。

\subsection{clipped surrogate objective}
PPO 的核心目标写成
\[
\tilde J(\theta')\approx \sum_{t,i}
\min\Bigl(
\rho_{t,i}\hat A_{t,i},
\text{clip}(\rho_{t,i},1-\epsilon,1+\epsilon)\hat A_{t,i}
\Bigr),
\]
其中
\[
\rho_{t,i}=\frac{\pi_{\theta'}(a_{t,i}\mid s_{t,i})}{\pi_\theta(a_{t,i}\mid s_{t,i})}.
\]
与单纯 importance weighting 相比，PPO 在这里增加了 clip，硬性限制新旧策略比值的变化范围。这样一来，若策略更新想走得过远，目标函数会自动截断收益，防止优化器因为局部看似有利就大步偏航。

\begin{figure}[H]
\centering
\includegraphics[width=0.96\textwidth]{pdf_pages/page-13.png}
\caption{PPO 的训练流程、clip 目标与代表性超参数}
\end{figure}

图中同时给出了一个重要实践信号：PPO 虽然允许在同一批数据上做多次梯度更新，但它依旧不是 fully off-policy。它追求的是“稳定的局部多次利用”，而不是无限制复用历史经验。

\subsection{GAE：用可调 horizon 估计 advantage}
PPO 还常与 generalized advantage estimation 配套。其基本思想是先得到一族不同步长的 advantage 估计
\[
\hat A_n^\pi(s_t,a_t)=\sum_{t'=t}^{t+n}\gamma^{t'-t}r(s_{t'},a_{t'})
-\hat V_\phi^\pi(s_t)+\gamma^n \hat V_\phi^\pi(s_{t+n}),
\]
再做加权和
\[
\hat A_{\text{GAE}}^\pi(s_t,a_t)=\sum_{n=1}^{\infty}w_n \hat A_n^\pi(s_t,a_t),
\qquad
w_n\propto \lambda^{n-1}.
\]
直觉上，较短 horizon 方差小、偏差大；较长 horizon 偏差小、方差大。GAE 用参数 \(\lambda\) 在两者之间平衡，使 advantage 更平滑、更适合作为 PPO 更新信号。

\begin{importantbox}{PPO 的本质}
PPO 不是把 RL 变简单了，而是把“多次更新会失控”这件事直接写进目标函数里约束住，因此成为一个稳定、常用、对实现友好的 actor-critic 算法。
\end{importantbox}

\subsection{本章小结}
PPO 通过 clip objective 和 GAE 两个关键组件，把“重要性采样修正”提升为真正好用的算法。它仍然基于 actor-critic，只是对目标构造和优势估计做了更强的工程化稳定处理。

\section{从局部 off-policy 到 fully off-policy：replay buffer}
PPO 只是在同一批数据上多更新几次，但本讲更进一步追问：能不能把所有过去的 trial-and-error 数据都留下来反复用？

\begin{figure}[H]
\centering
\includegraphics[width=0.96\textwidth]{pdf_pages/page-15.png}
\caption{replay buffer 版本的 off-policy actor-critic：每次都从历史经验中采样 minibatch}
\end{figure}

这就是 replay buffer 思路。算法不再只依赖当前 batch，而是维护一个经验池 \(\mathcal R\)，其中每条 transition 形如 \((s,a,s',r)\)。训练时从 \(\mathcal R\) 里抽样小批量，再执行参数更新。

它带来巨大样本效率优势，因为一段昂贵采样得到的数据，后面能被反复利用；但也立即引出两个理论和实现问题。

\subsection{问题一：为什么不能继续直接拟合状态价值}
若 replay buffer 中的数据来自大量旧策略，那么“状态 \(s\) 的未来回报”已经与旧策略动作分布纠缠在一起。直接拟合 \(V^\pi(s)\) 时，标签到底对应哪个策略，变得含糊不清。

因此讲者提出更合适的思路：改学 \(Q^\pi(s,a)\)。因为数据集中本来就存了动作 \(a\)，把动作当输入，就不必假装这些旧 transition 是由当前策略动作分布生成的。

\subsection{问题二：策略更新时应该信谁的动作}
即便有了 \(Q(s,a)\)，策略更新也不能直接拿 replay buffer 里旧动作照抄。因为旧策略的动作常常已经不是当前最优选择。策略梯度中的动作最好仍由当前 actor 采样，而 value estimator 则借助 replay buffer 提高样本利用率。

\begin{figure}[H]
\centering
\includegraphics[width=0.96\textwidth]{pdf_pages/page-20.png}
\caption{在 replay buffer 版本中，用 \(Q\) 替代 \(A\)，并让当前策略动作与历史动作分工}
\end{figure}

这张图非常关键。它说明 fully off-policy actor-critic 并不是“历史经验一把梭”，而是把 actor 与 critic 在数据使用上的角色重新分开：critic 更依赖 buffer，actor 更依赖当前策略采样或当前策略分布。

\subsection{本章小结}
replay buffer 让 actor-critic 的样本效率大幅提升，但也迫使我们放弃一些只在 on-policy 下自然成立的简化公式。正是在这里，课程开始通向 SAC 一类 practical off-policy algorithm。

\section{SAC 与算法取舍：数据效率、稳定性与应用场景}
本讲后段用 SAC 作为代表 practical algorithm。SAC 的细节将在后续继续展开，但这里已经给出其定位：它是 fully off-policy、基于 replay buffer 的 actor-critic 代表，在数据效率上通常明显优于 PPO。

\subsection{SAC 的高层图景}
SAC 继续沿用 replay buffer 与 \(Q\)-function 学习框架，同时把最大熵目标引入策略优化。直观地说，它不仅让 actor 追求高 \(Q\) 值，还鼓励策略保留一定随机性，以改进探索与鲁棒性。

从本讲范围看，更重要的不是推完整公式，而是理解 SAC 与 PPO 的分工：
\begin{itemize}
\item PPO 更稳定、实现更顺手、在大规模 on-policy 采样条件下很实用。
\item SAC 更省数据，但对超参数、目标构造和训练稳定性更敏感。
\end{itemize}

\begin{knowledgebox}{课程中的工程判断}
如果你最在乎稳定性和易用性，PPO 往往是默认起点；如果你最在乎数据效率，尤其真实机器人交互昂贵，SAC 常常更有吸引力。
\end{knowledgebox}

\subsection{为什么 RL for language models 常提 PPO}
讲者在结尾特别点出：PPO 是 language model RL 中常见的选择。原因并不神秘，而是因为 RLHF 这类流程天然更重视训练稳定性、约束更新幅度、避免策略突然崩坏。这和 PPO 的设计目标高度一致。

\subsection{本章小结}
本节给出 practical RL 的算法分工图谱：PPO 代表“稳定但不极致省样本”的 actor-critic，SAC 代表“更高样本效率但更难调”的 fully off-policy actor-critic。

\section{总结与延伸}
本讲真正完成的，不只是从 on-policy 到 off-policy 的字面变化，而是一次算法设计哲学的升级。上一讲的问题是“如何估值”；这一讲的问题变成“如何在估值的同时最大化复用数据，却不把训练搞崩”。

围绕这个目标，课程给出了两条主线。第一条主线是 PPO：接受数据仍然主要来自当前策略，但允许小心地多次复用同一批样本，并通过 clip 和 GAE 把不稳定性控制住。第二条主线是 replay buffer + \(Q\)-learning 式 critic：进一步拥抱历史数据，把数据效率推得更高，并最终导向 SAC。

这也解释了为什么 practical RL 没有单一最佳算法。不同任务在样本成本、动作空间、稳定性要求、工程复杂度上约束不同，导致 PPO 与 SAC 在工业和研究中长期并存。理解这点，比死记某一个目标函数更重要。

\end{document}
